% GENERATED FILE. Edit canonical lessons, then run npm run generate:books.
% canonical-source-hash: fnv1a64:24afb66a86f330f3
% canonical-lessons: GU-C430-trajvu, GU-C430-man, GU-C430-padiku, GU-C430-gharak, GU-C430-jokhvu

\chapter{Scales and Packets}
\label{ch:gu-scales-and-packets}
\hlchaptermodality{430}

\begin{chapteropening}
\textbf{By the end of this chapter:} \emph{I can buy loose goods by weight.}

Complete the last lesson of chapter 430: I can buy loose goods by weight.
\end{chapteropening}

\section[\texorpdfstring{\gu{ત્રાજવું} (trājvũ)}{ત્રાજવું (trājvũ)}]{\gu{ત્રાજવું} (trājvũ) — a pair of scales}

\label{lesson:GU-C430-trajvu}



\begin{quote}
Before the new one: say the Gujarati for to be robbed, then the Gujarati for to be affordable.
\end{quote}

\subsection*{You'll want to know: \gu{ત્રાજવું}}
\textbf{\gu{ત્રાજવું}} — \emph{trājvũ} — \textquotedblleft{}a pair of scales\textquotedblright{}.

\begin{quote}
\gu{દુકાનદાર} \gu{ત્રાજવામાં} \gu{ખાંડ} \gu{તોલે} \gu{છે}.
\end{quote}

\begin{quote}
\emph{dukāndār trājvāmā̃ khā̃ḍ tole che.}
\end{quote}

\begin{quote}
The shopkeeper weighs sugar on the scales.
\end{quote}

A neuter noun.

\begin{cousinweb}
\textbf{\gu{ત્રાજવું}} (\emph{trājvũ}) is usually traced to Persian \emph{tarāzū}, \textquotedblleft{}a balance, a pair of scales\textquotedblright{}.
\end{cousinweb}

\begin{grammarlens}[title={What you've built}]
One more word for getting around.
\end{grammarlens}

\subsection*{Your turn}
\begin{itemize}
\raggedright
  \item \emph{Say it:} \emph{trājvũ}
  \item \emph{Say it:} \emph{trājvũ}, once more
  \item \emph{Say it:} \emph{dukāndār trājvāmā̃ khā̃ḍ tole che.} — The shopkeeper weighs sugar on the scales.
  \item \emph{Recall:} say the Gujarati for to be robbed, then the Gujarati for to be affordable, then say \emph{trājvũ} again
\end{itemize}

\begin{tcolorbox}[breakable,colback=teal!4,colframe=teal!35!black,title={Before you move on}]
What does \gu{ત્રાજવું} mean? (\textquotedblleft{}A pair of scales\textquotedblright{}.) Say it once more.
\end{tcolorbox}
\section[\texorpdfstring{\gu{મણ} (maṇ)}{મણ (maṇ)}]{\gu{મણ} (maṇ) — a maund (an old measure of weight)}

\label{lesson:GU-C430-man}



\begin{quote}
Before the new one: say the Gujarati for to be affordable, then the Gujarati for a pair of scales.
\end{quote}

\subsection*{You'll want to know: \gu{મણ}}
\textbf{\gu{મણ}} — \emph{maṇ} — \textquotedblleft{}a maund (an old measure of weight)\textquotedblright{}.

\begin{quote}
\gu{વેપારીએ} \gu{બજારમાં} \gu{દસ} \gu{મણ} \gu{ઘઉં} \gu{વેચ્યા}.
\end{quote}

\begin{quote}
\emph{vepārīe bajārmā̃ das maṇ ghaũ vecyā.}
\end{quote}

\begin{quote}
The trader sold ten maunds of wheat in the market.
\end{quote}

A masculine noun.

\begin{cousinweb}
\textbf{\gu{મણ}} (\emph{maṇ}) is a tadbhava of Sanskrit \emph{maṇa}, \textquotedblleft{}a measure of grain\textquotedblright{}; CDIAL lists the Gujarati word, \textquotedblleft{}a maund\textquotedblright{} (CDIAL 9730). Farmers still sell grain by the maund.
\end{cousinweb}

\begin{grammarlens}[title={What you've built}]
One more word for getting around.
\end{grammarlens}

\subsection*{Your turn}
\begin{itemize}
\raggedright
  \item \emph{Say it:} \emph{maṇ}
  \item \emph{Say it:} \emph{maṇ}, once more
  \item \emph{Say it:} \emph{vepārīe bajārmā̃ das maṇ ghaũ vecyā.} — The trader sold ten maunds of wheat in the market.
  \item \emph{Recall:} say the Gujarati for to be affordable, then the Gujarati for a pair of scales, then say \emph{maṇ} again
\end{itemize}

\begin{tcolorbox}[breakable,colback=teal!4,colframe=teal!35!black,title={Before you move on}]
What does \gu{મણ} mean? (\textquotedblleft{}A maund (an old measure of weight)\textquotedblright{}.) Say it once more.
\end{tcolorbox}
\section[\texorpdfstring{\gu{પડીકું} (paḍīkũ)}{પડીકું (paḍīkũ)}]{\gu{પડીકું} (paḍīkũ) — a small packet (wrapped in paper)}

\label{lesson:GU-C430-padiku}



\begin{quote}
Before the new one: say the Gujarati for a pair of scales, then the Gujarati for a maund.
\end{quote}

\subsection*{You'll want to know: \gu{પડીકું}}
\textbf{\gu{પડીકું}} — \emph{paḍīkũ} — \textquotedblleft{}a small packet (wrapped in paper)\textquotedblright{}.

\begin{quote}
\gu{દુકાનદારે} \gu{ચાનું} \gu{પડીકું} \gu{વાળીને} \gu{આપ્યું}.
\end{quote}

\begin{quote}
\emph{dukāndāre cānũ paḍīkũ vāḷīne āpyũ.}
\end{quote}

\begin{quote}
The shopkeeper folded up a packet of tea and handed it over.
\end{quote}

A neuter noun.

\begin{cousinweb}
\textbf{\gu{પડીકું}} (\emph{paḍīkũ}) seems to belong with Sanskrit \emph{puṭa}, \textquotedblleft{}a fold, a pocket\textquotedblright{}, under which CDIAL lists a Gujarati word for a packet (CDIAL 8253).
\end{cousinweb}

\begin{grammarlens}[title={What you've built}]
One more word for getting around.
\end{grammarlens}

\subsection*{Your turn}
\begin{itemize}
\raggedright
  \item \emph{Say it:} \emph{paḍīkũ}
  \item \emph{Say it:} \emph{paḍīkũ}, once more
  \item \emph{Say it:} \emph{dukāndāre cānũ paḍīkũ vāḷīne āpyũ.} — The shopkeeper folded up a packet of tea and handed it over.
  \item \emph{Recall:} say the Gujarati for a pair of scales, then the Gujarati for a maund, then say \emph{paḍīkũ} again
\end{itemize}

\begin{tcolorbox}[breakable,colback=teal!4,colframe=teal!35!black,title={Before you move on}]
What does \gu{પડીકું} mean? (\textquotedblleft{}A small packet (wrapped in paper)\textquotedblright{}.) Say it once more.
\end{tcolorbox}
\section[\texorpdfstring{\gu{ઘરાક} (gharāk)}{ઘરાક (gharāk)}]{\gu{ઘરાક} (gharāk) — a customer (everyday); a regular buyer}

\label{lesson:GU-C430-gharak}



\begin{quote}
Before the new one: say the Gujarati for a maund, then the Gujarati for a small packet.
\end{quote}

\subsection*{You'll want to know: \gu{ઘરાક}}
\textbf{\gu{ઘરાક}} — \emph{gharāk} — \textquotedblleft{}a customer (everyday); a regular buyer\textquotedblright{}.

\begin{quote}
\gu{આ} \gu{દુકાનના} \gu{ઘરાક} \gu{વરસોથી} \gu{અહીં} \gu{આવે} \gu{છે}.
\end{quote}

\begin{quote}
\emph{ā dukānnā gharāk varsothī ahī̃ āve che.}
\end{quote}

\begin{quote}
This shop's customers have come here for years.
\end{quote}

A masculine noun.

\begin{cousinweb}
\textbf{\gu{ઘરાક}} (\emph{gharāk}) is usually traced to Sanskrit \emph{grāhaka}, \textquotedblleft{}one who takes\textquotedblright{} (Monier-Williams), worn down in speech; \gu{ગ્રાહક} keeps the learned form of the same word.
\end{cousinweb}

\begin{grammarlens}[title={What you've built}]
One more word for getting around.
\end{grammarlens}

\subsection*{Your turn}
\begin{itemize}
\raggedright
  \item \emph{Say it:} \emph{gharāk}
  \item \emph{Say it:} \emph{gharāk}, once more
  \item \emph{Say it:} \emph{ā dukānnā gharāk varsothī ahī̃ āve che.} — This shop's customers have come here for years.
  \item \emph{Recall:} say the Gujarati for a maund, then the Gujarati for a small packet, then say \emph{gharāk} again
\end{itemize}

\begin{tcolorbox}[breakable,colback=teal!4,colframe=teal!35!black,title={Before you move on}]
What does \gu{ઘરાક} mean? (\textquotedblleft{}A customer (everyday); a regular buyer\textquotedblright{}.) Say it once more.
\end{tcolorbox}
\section[\texorpdfstring{\gu{જોખવું} (jokhvũ)}{જોખવું (jokhvũ)}]{\gu{જોખવું} (jokhvũ) — to weigh out (goods)}

\label{lesson:GU-C430-jokhvu}



\begin{quote}
Before the new one: say the Gujarati for a small packet, then the Gujarati for a regular buyer.
\end{quote}

\subsection*{You'll want to know: \gu{જોખવું}}
\textbf{\gu{જોખવું}} — \emph{jokhvũ} — \textquotedblleft{}to weigh out (goods)\textquotedblright{}.

\begin{quote}
\gu{દુકાનદારે} \gu{બે} \gu{કિલો} \gu{ચોખા} \gu{જોખ્યા}.
\end{quote}

\begin{quote}
\emph{dukāndāre be kilo cokhā jokhyā.}
\end{quote}

\begin{quote}
The shopkeeper weighed out two kilos of rice.
\end{quote}

Past, with \gu{તેણે} (he or she …): \gu{તેણે} \gu{ખાંડ} \gu{જોખી}, \emph{teṇe khā̃ḍ jokhī}.

\begin{cousinweb}
\textbf{\gu{જોખવું}} (\emph{jokhvũ}) goes back to \emph{*yōkṣati}, a form CDIAL reconstructs; CDIAL lists the Gujarati verb, \textquotedblleft{}to examine, to weigh\textquotedblright{} (CDIAL 10525).
\end{cousinweb}

\begin{grammarlens}[title={What you've built}]
One more word for getting around.
\end{grammarlens}

\subsection*{Your turn}
\begin{itemize}
\raggedright
  \item \emph{Say it:} \emph{jokhvũ}
  \item \emph{Say it:} \emph{jokhvũ}, once more
  \item \emph{Say it:} \emph{dukāndāre be kilo cokhā jokhyā.} — The shopkeeper weighed out two kilos of rice.
  \item \emph{Recall:} say the Gujarati for a small packet, then the Gujarati for a regular buyer, then say \emph{jokhvũ} again
\end{itemize}

\begin{tcolorbox}[breakable,colback=teal!4,colframe=teal!35!black,title={Before you move on}]
What does \gu{જોખવું} mean? (\textquotedblleft{}To weigh out (goods)\textquotedblright{}.) Say it once more.
\end{tcolorbox}
